% TOP_PROBLEM: {"id": 3097, "title": "Fat 4-polytopes", "category_id": 6, "category": {"id": 6, "name": "geometry", "display_name": "Geometry", "description": "Euclidean and non-Euclidean geometry, geometric structures.", "slug": "geometry", "order_index": 6, "created_at": "2026-07-31T15:26:25.671Z"}, "status": "open", "problem_number": "OPG-316", "set_id": 12}
% FIRST_POSTED: 2026-10-01
% ATTEMPT_STATUS: unresolved
% COMPLETION: 5
% MODEL: GPT 6 Astra Ultra
% UnsolvedMath ID 3097; source catalogue code OPG-316
% !TeX program = lualatex
% Research attempt, 2026-10-01. Canonical statement preserved verbatim.
\documentclass[11pt,a4paper]{article}
\usepackage[margin=25mm]{geometry}
\usepackage{fontspec}
\setmainfont{lmroman10-regular.otf}[BoldFont=lmroman10-bold.otf,
 ItalicFont=lmroman10-italic.otf,BoldItalicFont=lmroman10-bolditalic.otf]
\usepackage{amsmath,amssymb,mathtools}
\usepackage{unicode-math}
\setmathfont{latinmodern-math.otf}
\AtBeginDocument{\let\setminus\smallsetminus}
\usepackage{xurl,hyperref}
\hypersetup{colorlinks=true,urlcolor=blue}
\setcounter{secnumdepth}{0}
\setlength{\parindent}{0pt}
\setlength{\parskip}{5pt}
\setlength{\emergencystretch}{3em}
\begin{document}
\section*{Fat 4-polytopes}\label{3097}
{\small UnsolvedMath ID 3097 · OPG-316 · Geometry · OpenGarden}
\subsection{Definitions and mathematical statement}
A convex $4$-polytope is the convex hull $P$ of finitely many points in
$\mathbb R^4$, with affine span all of $\mathbb R^4$. Its proper nonempty
faces are the intersections with supporting hyperplanes: a supporting
hyperplane meets $P$ and has $P$ entirely in one closed half-space.
For $0\le i\le3$, let $f_i(P)$ count its $i$-dimensional faces.
In particular, $f_0,f_1,f_2,f_3$ count vertices, edges, polygonal
$2$-faces and $3$-dimensional facets, respectively.

Use the precise normalization in the catalogue:
\[
 \operatorname{fat}(P)=\frac{f_1(P)+f_2(P)}{f_0(P)+f_3(P)}.
\]
The denominator is positive. The question asks whether there is a
finite real constant $C$ such that $\operatorname{fat}(P)\le C$ for
every convex $4$-polytope $P$.
Equivalently, is the supremum of this ratio over all convex
$4$-polytopes finite?

A negative answer must supply polytopes whose ratios grow without bound:
for every $R>0$, some $P$ must have $f_1(P)+f_2(P)>R(f_0(P)+f_3(P))$.
A positive answer may use any absolute constant independent of the
number of vertices, number of facets, coordinates and combinatorial type.
The question concerns boundedness, not the exact optimal constant.

No simplicial or simple hypothesis is imposed. Only convex polytopes are
included; abstract cell decompositions or nonconvex realizations do not
by themselves answer it. Some literature modifies the ratio by subtracting
constants from numerator and denominator; this definition retains the
unshifted ratio appearing in this UnsolvedMath record.
\subsection{Short English statement}
Is the ratio of edges plus two-dimensional faces to vertices plus facets bounded over all convex four-dimensional polytopes?
\subsection{Research attempt}
\paragraph{Scope and literature check.}
Eppstein, Kuperberg and Ziegler [A1] distinguish convex $4$-polytopes from
abstract sphere decompositions and construct convex examples with fatness
above five. Their unbounded examples in a broader topological class cannot
be treated as convex examples. We use exactly the unshifted ratio from the
canonical definition. This attempt quantifies why simple polytopes, their
duals, and elementary products cannot establish unboundedness. The primary
abstract and author manuscript were checked on 1 October 2026.

\paragraph{A vertex-degree bound.}
Write $v=f_0$, $e=f_1$, $r=f_2$, $f=f_3$. The boundary of a convex
$4$-polytope is a cell decomposition of $S^3$, so its Euler relation is
$v-e+r-f=0$. Hence
\[
 \operatorname{fat}(P)=\frac{2e+f-v}{v+f}.
\]
If every vertex has degree at most $D$, the handshake identity $2e\le Dv$
gives
\[
 \operatorname{fat}(P)\le\frac{(D-1)v+f}{v+f}\le D-1
\]
for $D\ge2$. This estimate uses no simplicity assumption. By applying it
to the polar polytope, the same bound holds if every facet has at most $D$
ridges: polar duality exchanges $(v,e,r,f)$ with $(f,r,e,v)$ and leaves
fatness unchanged. Thus a family with unbounded fatness must have both
unbounded maximum vertex degree and unbounded maximum number of ridges per
facet. A bound on either one alone blocks the proposed construction.

\paragraph{The exact simple and simplicial barriers.}
For a simple $4$-polytope every vertex has degree four, so $e=2v$ and
$r=v+f$. Consequently
\[
 \operatorname{fat}(P)=\frac{3v+f}{v+f}=3-\frac{2f}{v+f}<3.
\]
For a simplicial $4$-polytope the dual is simple, giving the same strict bound.
It follows that merely adding vertices to simple polytopes, or taking many
simplices in a convex simplicial construction, cannot drive this particular
ratio beyond three. Triangulating facets also changes the face counts and
need not preserve the original polytope's fatness; a triangulated sphere is
not a substitute for a convex realization with those faces.

\paragraph{An explicit family approaching the barrier.}
Let $Q_m,Q_n$ be convex polygons with $m,n\ge3$ vertices and set
$P=Q_m\times Q_n$. Exposing a face by a linear functional on the product
exposes a face in each factor. Counting all products of faces of total
dimension $0,1,2,3$ gives
\[
 (v,e,r,f)=(mn,2mn,mn+m+n,m+n).
\]
For instance, the $2$-faces comprise $mn$ edge-products, $n$ copies of the
first polygon, and $m$ copies of the second. Therefore
\[
 \operatorname{fat}(P)=3-\frac{2(m+n)}{mn+m+n}\longrightarrow3
 \quad\text{as }m=n\longrightarrow\infty.
\]
The family has arbitrarily many vertices and edges but remains bounded.
For two triangles its vector is $(9,18,15,6)$ and fatness is $33/15$;
for two squares it is $(16,32,24,8)$ and fatness is $56/24$.

\paragraph{Verification and first remaining gap.}
The counts satisfy Euler's relation identically and are unchanged in their
fatness ratio under duality. The strict inequality below three depends on
$f>0$, as required for any full-dimensional bounded polytope. The general
bound is on maximum degree, not average degree silently rounded to a fixed
constant. Since prior work already exceeds five, the restricted barrier
must not be reported as a universal upper bound. A successful unbounded
construction needs more complicated vertex and facet incidences together
with convex realizability. We have neither such a construction nor an
inequality bounding that unrestricted incidence complexity.

\paragraph{Final status.}
\textbf{UNRESOLVED.} The result is a proved obstruction to several elementary
families, including exact limiting values for polygon products. It does not
answer boundedness for arbitrary convex $4$-polytopes. No novelty claim is
made for the identities. Completion estimate: 5\%.

\subsection{Sources}
\begin{itemize}
\item[S1] ULAM.AI and the UnsolvedMath Contributors, supplied problems.json, record 3097 (OPG-316), OpenGarden collection. Catalogue identity and source transcription; CC BY 4.0.
\url{https://huggingface.co/datasets/ulamai/UnsolvedMath}.
\item[S2] Open Problem Garden, Fat 4-polytopes. Original problem and discussion; the mathematical formulation below is expanded from the supplied source record.
\url{http://www.openproblemgarden.org/op/fat_4_polytopes}.
\item[A1] D. Eppstein, G. Kuperberg and G. M. Ziegler, \emph{Fat 4-polytopes and fatter 3-spheres}, arXiv:math/0204007. \url{https://arxiv.org/abs/math/0204007}. Author manuscript: \url{https://www.mi.fu-berlin.de/math/groups/discgeom/ziegler/Preprintfiles/080PREPRINT.pdf}.
\end{itemize}
\end{document}
